% ====================================================================
% graphite_ica_chapter1_v0.1.tex
% Chapter 1 Rebuild (from scratch) — Anode Fit Project
% Date: 2026-05-27
% Author: Claude (Project_Anode_Fit Claude 측)
% ====================================================================
%
% CHARTER COMPLIANCE
% ------------------
% This file is governed by the Chapter 1 Rebuild Charter:
%   Claude/results/PHASE_E0_foundation_reset_charter_RESULT.md
%
% Anti-pattern list (Charter §2 — MUST NOT appear in definitional eqs):
%   AP1 max(ΔG_eff, 0)   AP2 min(., k_max)        AP3 logistic
%   AP4 erf as definitional   AP5 discrete N_p as primary state
%   AP6 logistic in rechecked2   AP7 softplus as canonical   AP8 N_p retained
%
% Forbidden-pattern list (Charter §4):
%   FB1 max/Heaviside, FB2 min/clip, FB3 sigmoid/tanh, FB4 erf/erfc,
%   FB5 discrete N_p primary, FB6 softplus without smooth-limit,
%   FB7 sign/clip/piecewise, FB8 silent μ discretization,
%   FB9 silent integral→sum, FB10 single-state simplification
%
% Smooth-limit consistency rule (Charter §5):
%   Every regularization parameter ε>0 has an explicit ε→0 form
%   recovering a Charter-compliant continuous expression.
%
% Audit Dim #11 (Charter §6) — system-fidelity:
%   Pass 1: enumerate every max/min/clip/step/Heaviside/sign/softplus/
%           logistic/erf/sigmoid/tanh/relu occurrence.
%   Pass 2: classify FAIL-definitional / WARN-regularization /
%           WARN-empirical / OK-derived.
%   Pass 3: confirm 0 remaining FAIL.
%
% Spine (Phase E1 §A.3):
%   (1) Q_ext = Q_cell q
%   (2) Q_cell q = Q_bg(V_n, T) + ∫dμ ρ(μ; q, T) n(μ)                    [Eq. 48]
%   (3) V_n = root of (2)                                                  [Loop A]
%   (4) V_{n,app} = V_n + s_I |I| R_n
%   (5a) ρ(μ; q, t) = ρ_0 + ∫dμ' K_n S_R [ρ_eq - ρ]                      [Eq. 49]
%   (5b) ICA = Q_cell / (dV_{n,app}/dq)
%   (5c) DVA = (1/Q_cell)(dV_{n,app}/dq)
% ====================================================================

\documentclass[11pt,a4paper]{article}
\usepackage[margin=22mm]{geometry}
\usepackage{kotex}
\usepackage{amsmath,amssymb,mathtools,bm}
\usepackage{booktabs,longtable,array,tabularx,multirow}
\usepackage{enumitem}
\usepackage{hyperref}
\usepackage{xcolor}
\usepackage{titlesec}
\usepackage{setspace}
\usepackage{fancyhdr}
\usepackage{lastpage}
\usepackage{caption}
\usepackage{float}
\usepackage{amsthm}
\usepackage{mdframed}

\hypersetup{colorlinks=true, linkcolor=blue, urlcolor=blue, citecolor=blue,
            pdftitle={Chapter 1 Rebuild — Continuous Chemical Potential ICA/DVA},
            pdfauthor={Project_Anode_Fit Claude}}
\setstretch{1.14}
\setlength{\parskip}{0.42em}
\setlength{\parindent}{0pt}
\setlength{\headheight}{14pt}
\pagestyle{fancy}
\fancyhf{}
\lhead{Chapter 1 (Rebuild v0.1) — 연속 화학 퍼텐셜 기반 ICA/DVA}
\rhead{\thepage/\pageref{LastPage}}
\renewcommand{\headrulewidth}{0.4pt}

\newtheorem*{assumption}{가정}
\newtheorem*{note}{설명}
\newtheorem*{charter}{Charter compliance}

% Preserved macros (ver5 + rechecked2 common subset)
\newcommand{\dd}{\mathrm{d}}
\newcommand{\ee}{\mathrm{e}}
\newcommand{\OCV}{\mathrm{OCV}}
\newcommand{\app}{\mathrm{app}}
\newcommand{\drive}{\mathrm{drive}}
\newcommand{\eff}{\mathrm{eff}}
\newcommand{\eq}{\mathrm{eq}}
\newcommand{\tot}{\mathrm{tot}}
\newcommand{\bg}{\mathrm{bg}}
\newcommand{\cell}{\mathrm{cell}}
\newcommand{\ext}{\mathrm{ext}}
\newcommand{\Ah}{\mathrm{Ah}}
\newcommand{\Cunit}{\mathrm{C}}

% New macros for continuous-form variables (Phase E0 §3.1, Phase E1 §B)
\newcommand{\mucont}{\mu}
\newcommand{\rhomu}{\rho(\mucont;\, q, T)}
\newcommand{\rhoeq}{\rho_{\eq}(\mucont;\, q, T)}
\newcommand{\rhomut}{\rho(\mucont;\, q, t)}
\newcommand{\Sreact}{S_R(\mucont;\, T)}
\newcommand{\Kncross}{K_n(\mucont, \mucont';\, q, T)}
\newcommand{\nmu}{n(\mucont)}

\title{\Large \textbf{리튬이온전지 흑연 음극의 연속 화학 퍼텐셜 기반 ICA/DVA 동역학:
       Chapter 1 (Rebuild v0.1)}}
\author{Project\_Anode\_Fit Claude 측}
\date{2026-05-27}

\begin{document}
\maketitle

\begin{center}
\begin{minipage}{0.92\textwidth}
\textbf{초록.} 이 문서는 리튬이온전지 흑연 음극의 ICA($dQ/dV$)와 DVA($dV/dQ$)를
\textbf{연속 화학 퍼텐셜 분포} $\rhomu$와 \textbf{연속 반응성 커널} $\Sreact$로
모델링하는 Chapter 1의 처음부터 다시 쓴 판본이다. 자기 참조 적분식의 해는 사용자
박사학위 연구에서 비롯된 \textbf{Fredholm 적분 방정식 제2종의 비율 치환법}
(Refs.~\cite{lee2011propagator,son2013accurate}, JCP 2017)을 적용한다. 본 Chapter 1
은 기존 \texttt{ver5} 와 \texttt{ver1\_rechecked2} 의 본문을 수정하지 않고 참조
자료로만 사용하며, 사용자가 지적한 ``적분을 모 아니면 도 즉 스텝 펑션의 형태로
가정'' 한 근본 결함을 회피하기 위해 진행률 변수 $\xi_j$를 본 spine 의 primary
state 에서 배제하고 $\rhomu$ 로부터의 \emph{유도} 양으로 격하한다.
\end{minipage}
\end{center}

\tableofcontents
\newpage

% ====================================================================
\section{문서의 목적과 원칙}
% ====================================================================

\subsection*{1.1 문서의 목적}

본 문서는 흑연 음극의 staging transition 거동을 외부 관측량 $dQ/dV$, $dV/dQ$에
연결하기 위한 Chapter 1의 \textbf{재작성판 v0.1} 이다. 모델의 중심 상태 변수로
\textbf{연속 화학 퍼텐셜 분포} $\rhomu$ 를 두며, 반응성을 \textbf{연속 커널}
$\Sreact$ 로 표현하고, 그로부터 발생하는 자기 참조 적분 방정식의 해를
\textbf{Fredholm 적분 방정식 제2종의 비율 치환법} 으로 닫힌 형태로 구한다.

\subsection*{1.2 4 가지 원칙}

본 Chapter 1 의 모든 정의식과 유도식은 다음 4 가지 원칙을 따른다.

\begin{enumerate}[label=\textbf{(P\arabic*)}]
\item \textbf{연속성 우선.} 본 spine 의 모든 정의식은 화학 퍼텐셜 좌표 $\mucont$
      에 대해 \emph{연속} 이다. 이산적 $N_p$ 전이 분해, logistic/erf/sigmoid/softplus
      을 사용한 평형 분포 정의, $\max(\cdot, 0)$ 또는 $\min(\cdot, k_{\max})$
      형태의 단단한 절단을 정의식에 사용하지 않는다. 수치 안정성을 위한 정규화는
      허용하되, 명시적 $\varepsilon \to 0$ smooth-limit 절을 함께 둔다.

\item \textbf{전하 보존식 중심.} 음극 내부 전위 $V_n$ 은 외부 함수 lookup 으로
      주어지지 않고, 본 Chapter 의 중심식인 전하 보존식 (\S\ref{sec:charge_balance})
      의 \emph{implicit 해} 로 정의된다. $V_n$ 은 분포 $\rhomu$ 의 화학 퍼텐셜
      공간 경계값으로 해석된다.

\item \textbf{자기 참조 적분식의 닫힌 해.} 본 모델에서 발생하는 자기 참조 적분
      방정식 (Fredholm 제2종) 은 사용자 본인의 박사학위 연구
      \cite{lee2011propagator, son2013accurate} 에서 제시된 \emph{비율 치환법}
      으로 닫힌 형태의 해석해를 얻는다. 단순 반복 해법이 발산하는 영역까지 적용
      가능한 방법이다.

\item \textbf{이전 판본 보존.} 기존 \texttt{ver5} (Chapter 1 \(\sim\) Chapter 5
      통합 적층본) 와 \texttt{ver1\_rechecked2} (Chapter 1 재작성 시도판) 의 본문은
      \emph{수정하지 않는다}. 본 Chapter 1 v0.1 은 그 둘을 참조 자료 (특히
      anti-pattern 의 근거 위치 및 평형 한계의 비교) 로만 사용한다.
\end{enumerate}

\subsection*{1.3 적용 범위}

본 Chapter 1 의 적용 범위는 다음과 같다.

\begin{itemize}
\item 흑연 음극의 저율 (OCV 근사) 부터 유한 C-rate 까지의 $dQ/dV$, $dV/dQ$ 해석
\item OCV 기준 또는 0.2C 기준 전위 곡선과의 연결
\item 등온 조건이 기본이며, 온도 의존성은 $\rhomu, \rhoeq, \Sreact, R_n$ 의 $T$
      파라미터로 inline 반영 (별도 ``온도 효과'' 단원 없음)
\item 충전 branch 의 히스테리시스는 본 Chapter 1 범위 외 — Chapter 5 (재작성 시
      별 로드맵) 로 이관
\end{itemize}

\subsection*{1.4 기존 판본과의 관계}

본 Chapter 1 v0.1 은 \texttt{ver5} 와 \texttt{ver1\_rechecked2} 의 \emph{디벨롭이
아니라} 처음부터의 \emph{재작성} 이다.

\begin{itemize}
\item \texttt{ver5} 의 본질적 결함: 진행률 $\xi_j$ 의 평형값 $\xi_{j,\eq}$ 를
      logistic 또는 erf 로 정의하여 staging 을 본질적으로 step function 의 합으로
      가정하고, 활성화 장벽을 $\max(\Delta G_{\eff,j}, 0)$ 또는
      $\min(k_j, k_{\max})$ 로 절단하는 등, 사용자가 지적한
      ``적분을 모 아니면 도 즉 스텝 펑션 형태'' 의 가정이 spine 자체에 박혀 있다.
\item \texttt{ver1\_rechecked2} 의 부분 개선과 잔존 문제: $V_n$ 을 전하 보존식의
      implicit 해로 격상하고 softplus 정규화를 도입했으나, $\xi_{j,\eq}$ logistic
      과 $N_p$ 이산 분해는 그대로 유지되어 본질이 변경되지 않았다. 10 회 이상의
      검토 과정에서 ``특정 변수가 되먹이는 수식'' 등 잔존 문제가 반복 검출되어
      디벨롭으로는 해결 불가하다고 사용자가 판단한 상태이다.
\item 본 Chapter 1 v0.1 의 출발점: 위 두 판본의 결함이 \emph{spine 자체} 에
      있으므로, spine 을 처음부터 새로 설계한다. 진행률 $\xi_j$ 는 primary
      state 에서 폐기되고, 연속 분포 $\rhomu$ 가 그 자리를 대신한다.
\end{itemize}

\subsection*{1.5 본 Chapter 1 의 spine 요약}

본 spine 은 5 개 성분으로 구성된다 (자세한 정의는 후속 \S 들에서).

\begin{equation}
\boxed{
\begin{aligned}
\text{(1)} \quad & Q_{\ext} = Q_{\cell}\, q \\
\text{(2)} \quad & Q_{\cell}\, q = Q_{\bg}(V_n, T) +
                   \int_{\mucont_{\min}}^{\mucont_{\max}} \dd\mucont \;
                   \rhomu \cdot \nmu \\
\text{(3)} \quad & V_n = \text{(2) 의 } V_n \text{ 에 대한 implicit 해} \\
\text{(4)} \quad & V_{n,\app} = V_n + s_I |I|\, R_n(q, T, |I|) \\
\text{(5a)} \quad & \rhomut = \rho_0(\mucont; q, t) +
                   \int_{\mucont_{\min}}^{\mucont_{\max}} \dd\mucont'\;
                   \Kncross \cdot S_R(\mucont'; T) \cdot
                   \bigl[\rhoeq - \rho(\mucont'; q, t)\bigr] \\
\text{(5b)} \quad & \frac{\dd Q_{\ext}}{\dd V_{n,\app}}
                   = \frac{Q_{\cell}}{\dd V_{n,\app}/\dd q} \\
\text{(5c)} \quad & \frac{\dd V_{n,\app}}{\dd Q_{\ext}}
                   = \frac{1}{Q_{\cell}}\, \frac{\dd V_{n,\app}}{\dd q}
\end{aligned}
}
\label{eq:spine_summary}
\end{equation}

\begin{note}
본 spine 의 (5a) 는 $\rhomut$ 에 대한 \textbf{Fredholm 적분 방정식 제2종} 이다.
$q$ 슬라이스 별로 보았을 때 그 구조는 사용자 본인의 JCP 2017 논문 식 (32) 과
1:1 대응되며, 비율 치환법으로 닫힌 형태의 해를 얻을 수 있다. 자세한 적용은
\S 17 (Phase E11 작성 단원) 에서 다룬다.
\end{note}

\begin{mdframed}
\textbf{이 문서는 진행률 변수 $\xi_j$ 를 primary state 로 사용하지 않는다.}
$\xi_j$ 는 연속 분포 $\rhomu$ 의 $j$ 번째 peak window 적분 으로부터 \emph{유도}
되는 양으로 격하된다. 이 격하는 사용자가 지적한 step function 가정의 본질적 회피
이며, 본 Charter (Phase E0 §3, §4) 의 핵심이다.
\end{mdframed}

% ====================================================================
\section{기호와 단위 컨벤션}
% ====================================================================

\subsection{기호 정의}
\label{sec:symbols}

본 Chapter 1 의 기호는 \textbf{canonical}, \textbf{derived}, \textbf{parameter}
3 종으로 분류된다. Canonical 기호는 spine 의 정의에 직접 등장하는 primary state
또는 입력 변수이고, derived 기호는 canonical 로부터 적분 또는 한계로 정의되는
보조량이며, parameter 는 외부에서 주어지거나 피팅 대상이 되는 상수 또는 함수
파라미터이다.

\subsubsection*{2.1.1 Canonical 기호 (17)}

\begin{longtable}{p{0.18\textwidth} p{0.16\textwidth} p{0.62\textwidth}}
\toprule
기호 & 단위 & 의미 \\
\midrule
\endhead
$\mucont$ & V (Li/Li$^+$ 기준) & 연속 화학 퍼텐셜 좌표. 본 spine 의 primary
    distribution axis. \\
$\rhomu$ & C/V & 화학 퍼텐셜 $\mucont$ 에서의 전하 밀도 (Li 점유량 밀도). 본
    spine 의 primary state variable. \\
$\Sreact$ & (단위 Phase E6 확정) & $\mucont$ 위치의 연속 반응성 커널. 이산 $k_j$
    를 대체한다. \\
$\Kncross$ & (단위 Phase E6 확정) & $\mucont$ 와 $\mucont'$ 사이의 연속 교차 커널
    (이산 $\rho_j(g)$ 를 대체). \\
$\rhoeq$ & C/V & 주어진 $(q, T)$ 에서의 평형 분포. $\rhomut$ 의 완화 표적이다. \\
$Q_{\bg}(V_n, T)$ & C & 잔류 chemical capacitance 함수. \texttt{ver1\_rechecked2}
    §4.1 의 역할을 보존한다. \\
$V_n$ & V & 음극 내부 전위. 식 \eqref{eq:spine_summary} (2) 의 $V_n$ 에 대한
    implicit 해. \\
$V_{n,\app}$ & V & 외부 관측 음극 전위. $V_n + s_I |I| R_n$. \\
$V_{n,\drive}$ & V & 반응성 커널의 구동 전위. reduced model 에서는 $V_{n,\app}$
    근사. \\
$R_n(q, T, |I|)$ & $\Omega$ & 음극 유효 저항 또는 유효 분극 계수. \\
$Q_{\cell}$ & C & 셀 기준 용량 (쿨롱 단위 고정). \\
$q$ & 무차원 & 방전 진행 좌표, $Q_{\ext}/Q_{\cell}$. \\
$I$ & A & 셀 전류 크기. \\
$T$ & K & 절대 온도. \\
$t$ & s & 시간. \\
$s_I,\, s_{\phi,j}$ & $\pm 1$ & 부호 컨벤션 (전류 방향, 상변이 방향). 충전 branch
    처리는 Chapter 5 로 이관. \\
$\nmu$ & 무차원 & 화학 퍼텐셜 단위 당 Li 점유 multiplier. 기본 선택 $\nmu = 1$
    (Phase E5 재검토 가능). \\
\bottomrule
\end{longtable}

\subsubsection*{2.1.2 Derived 기호 (legacy compatibility)}

다음 기호는 spine 의 canonical 변수로부터 적분 또는 한계로 \emph{유도} 되는 양이며,
주로 기존 \texttt{ver5}/\texttt{ver1\_rechecked2} 와의 호환성 또는 피팅 단계
(\S 13) 에서 사용된다. \textbf{본 spine 의 정의식에는 등장하지 않는다.}

\begin{longtable}{p{0.18\textwidth} p{0.62\textwidth} p{0.18\textwidth}}
\toprule
기호 & 유도식 & 최초 사용 \S \\
\midrule
\endhead
$Q_{\ext}$ & $Q_{\cell} \cdot q$ & §1 \\
$V_{n,\OCV}(q, T)$ & $\rho = \rhoeq$ 하에서의 $V_n$ (평형 분포의 경계값) & §11
    ($\rho^{\text{simple}}$ 구성) \\
$\xi_j$ & $(1/N_j) \int_{\mucont \in \text{peak}_j} \dd\mucont\, \rhomut$ & §10
    (피팅) \\
$\xi_{j,\eq}$ & $\rhomut \to \rhoeq$ 치환 & §10 \\
$k_j$ & $\Sreact|_{\mucont \approx U_j}$ & §10 \\
$U_j(T)$ & $\rhoeq$ 의 peak $j$ 중심 & §10 \\
$w_j(T)$ & $\rhoeq$ 의 peak $j$ FWHM & §10 \\
$Q_{j,\tot}$ & peak $j$ window 의 $q$ 진행 적분 & §10 \\
$\dd\xi_j/\dd q$ & $(1/N_j) \int_{\text{peak}_j} \dd\mucont\, \partial \rho / \partial q$ & §10 \\
$\mathcal A_j$ & $\Sreact$ 와 $V_{n,\drive}$ 의 결합으로부터 & §7 \\
$\Delta G_{a,j}$ & $\Sreact$ 의 Marcus 류 파라미터 & §7 \\
$\Delta G_{\eff,j}$ & $\mucont \approx U_j$ 에서 evaluated barrier & §7 \\
$\nu_j(T)$ & $\Sreact$ 의 prefactor at $\mucont \approx U_j$ & §7 \\
$\chi_j$ & 커널의 coupling coefficient & §7 \\
$\rho_j(g)$ (\texttt{ver5}) & $\Kncross$ 로 대체 & §7 \\
$g$ (\texttt{ver5}) & $\mucont, \mucont'$ 로 대체 & — \\
\bottomrule
\end{longtable}

\begin{charter}
위 derived 표의 모든 행은 \emph{유도} 양으로 명시되어 있다. 본 spine 의 정의식
(\S 4, \S 5, \S 6, \S 7) 에서는 derived 기호가 등장하지 않는다. 만약 후속 \S 에서
derived 기호가 정의식에 사용된다면, Audit Dim \#11 Pass 2 의 FAIL 조건 (Charter
\S 6.2) 에 해당한다.
\end{charter}

\subsection{단위 컨벤션}

본 Chapter 의 미분 방정식은 모두 SI 단위로 표기된다.

\begin{itemize}
\item $Q_{\cell}$ 은 \textbf{쿨롱 (C)} 단위로 고정한다. 셀 데이터가 Ah 단위로
      주어진 경우 다음 변환을 사용한다.
      \begin{equation}
      Q_{\cell}[\Cunit] = 3600 \cdot Q_{\cell}[\Ah].
      \label{eq:ah_to_c}
      \end{equation}
      이 변환을 생략하면 $\dd q/\dd t = |I|/Q_{\cell}$ 의 차원이 깨진다.
\item $\mucont$ 는 \textbf{V (Li/Li$^+$ 기준)} 단위이다. 외부 절대 전극 기준을
      바꾸는 경우 본 Chapter 의 모든 $\mucont$ 값에 동일한 평행 이동을 적용한다.
\item $\rhomu$ 는 \textbf{C/V} 단위이다. 즉
      $\int \dd\mucont\, \rhomu \cdot \nmu$ 의 차원이 \textbf{C} 가 되어 식
      \eqref{eq:spine_summary} (2) 의 양변 차원이 일치한다.
\item $\Sreact$ 와 $\Kncross$ 의 정확한 단위는 \S 7 의 커널 구성 단계 (Phase
      E6 작성) 에서 확정된다. 본 Phase E2 에서는 plausibly $\Sreact \sim 1/$s,
      $\Kncross \sim 1/$(V$\cdot$s) 로 표기하되, 최종 단위 정합성은 \S 7 본문에서
      재검토한다 (OI-E0-1).
\item $R_n$ 은 $\Omega$, $I$ 는 A, $T$ 는 K, $t$ 는 s 로 한다. $q$ 는 무차원이다.
\end{itemize}

\subsection{전류와 전위 부호 컨벤션}

방전 방향에서 음극 탈리튬화로 음극 전위가 증가하는 컨벤션을 사용한다. 전류 크기는
양수 $|I| \ge 0$ 로 두고, 방향은 부호 인자 $s_I$ 로 처리한다.

\begin{equation}
V_{n,\app} = V_n + s_I\, |I|\, R_n(q, T, |I|),
\qquad s_I = +1 \;\text{(방전 기본)}.
\label{eq:Vapp_def}
\end{equation}

상변이 방향 부호 $s_{\phi,j}$ 는 (legacy $\xi_j$ 기반 피팅 단계에서만) peak $j$
의 진행 방향과 전위 변화 방향의 결합 부호로 사용된다. \textbf{spine 의 정의식
에서는 $s_{\phi,j}$ 가 등장하지 않는다.}

충전 branch 의 처리 (양방향 $s_I$, 비대칭 $\rhoeq$ 등) 는 Chapter 5 (히스테리시스)
의 재작성 로드맵에서 별도 정의된다. 본 Chapter 1 v0.1 은 방전 branch 단일 처리만
다룬다.

\subsection*{2.4 본 \S 의 \S 1 spine 과의 연결}

본 \S 2 의 기호 정의는 \S 1 식 \eqref{eq:spine_summary} 의 모든 변수를 1:1 로
포괄한다. \S 1 의 spine (1) $\sim$ (5c) 에 등장하는 기호 $Q_{\ext}, Q_{\cell}, q,
Q_{\bg}, V_n, \rhomu, \nmu, V_{n,\app}, s_I, |I|, R_n, \rhomut, \rho_0, \Kncross,
\Sreact, \rhoeq$ 는 본 \S 2 의 canonical 표 \ref{sec:symbols} 에 모두 등재되어
있다. 따라서 후속 \S 의 본문은 \S 2 의 정의 위에서 자유롭게 전개될 수 있다.

% ====================================================================
\section{흑연 staging 과 연속 화학 퍼텐셜}
\label{sec:effective_transition}
% ====================================================================

\subsection{$N_p$ 이산 분해에서 연속 분포로의 전환}
\label{sec:np_to_continuous}

기존 \texttt{ver5} Chapter 1 \S 3 (line 137--150) 와 \texttt{ver1\_rechecked2}
\S 3 (line 109--116) 은 흑연 음극의 staging transition 을 $N_p$ 개 (보통 3 또는 4)
의 \emph{이산 effective transition} 으로 분해하고, 각 transition $j$ 에 대해 진행률
$\xi_j$ 를 primary state 로 두었다. 이 결정은 평형 진행률 $\xi_{j,\eq}$ 의 정의를
logistic 또는 erf 형태로 강제하고 (Charter \S 2 AP3, AP4, AP6), 결과적으로 spine
전체가 step function 의 합 형태로 환원되는 본질적 결함을 만들었다.

본 Chapter 1 v0.1 은 이 이산 분해를 \textbf{primary state 에서 제거} 하고, 흑연
staging 을 화학 퍼텐셜 좌표 $\mucont$ 에 대한 \textbf{연속 분포} $\rhomu$ 로
모델링한다. $N_p$ 개 peak 가 관측되는 사실은 $\rhomu$ 가 $\mucont$-공간에서
$N_p$ 개의 국소 집중 (peak) 을 갖는다는 형태로 자연스럽게 표현되며, 분리 가능 여부
는 peak 의 폭과 간격에 의해 결정되는 \emph{연속 모델의 결과} 이다.

\begin{assumption}
$\mucont$ 는 음극 내부 Li 원자의 (Li/Li$^+$ 기준) 전기화학 퍼텐셜로 해석한다.
$\rhomu \, \dd\mucont$ 는 $\mucont$ 주변 단위 화학 퍼텐셜 구간에서 흑연 호스트가
수용하는 Li 점유량의 전하 밀도 (단위 C/V) 이다. 이 해석은 사용자 검토 항목
(\textbf{DQ-G2}) 이며, 본 Chapter 1 의 모든 정의식은 이 해석 하에서 일관된다.
\end{assumption}

\subsection{$\rhomu$ 의 정의와 의미}
\label{sec:rho_definition}

연속 분포 $\rhomu$ 는 다음 의미를 갖는다.

\begin{equation}
\rhomu \, \dd\mucont \;=\;
\text{화학 퍼텐셜 구간 } [\mucont,\, \mucont + \dd\mucont] \text{ 에서의 Li
점유 전하 밀도}
\label{eq:rho_meaning}
\end{equation}

따라서 화학 퍼텐셜 구간 $[\mucont_a, \mucont_b]$ 에 누적된 Li 전하량은
\begin{equation}
Q_{[\mucont_a, \mucont_b]}(q, T)
\;=\;
\int_{\mucont_a}^{\mucont_b} \dd\mucont \; \rhomu \cdot \nmu
\label{eq:rho_window_integration}
\end{equation}
이며, 전체 화학 퍼텐셜 범위 $[\mucont_{\min}, \mucont_{\max}]$ 에서 누적된 Li
전하량은 식 \eqref{eq:spine_summary} (2) 의 우변 적분 항과 같다. 즉 본 spine 의
중심식 (전하 보존식, \S\ref{sec:charge_balance} 에서 정식 유도) 은 식
\eqref{eq:rho_window_integration} 을 전 범위에 적용한 형태이다.

$\rhomu$ 의 정의는 다음 3 가지 자유도를 갖는다.
\begin{enumerate}[label=(\alph*)]
\item $\mucont$-공간 분포의 형태 — Gaussian, Marcus 류, Cahn-Hilliard 류 또는
      그 혼합. 본 \S 3 에서는 형태를 특정하지 않으며, 형태의 선택은
      \S\ref{sec:rho_eq} (Phase E5) 에서 다룬다.
\item 분포의 $q$-의존성 — Li 채움이 진행될수록 어느 $\mucont$ 영역에 점유가
      증가하는지. 본 \S 3 에서는 정성적 설명만 두며, 정량은 \S\ref{sec:rho_eq}
      에서.
\item 분포의 $T$-의존성 — 온도가 peak 의 폭과 위치를 어떻게 바꾸는지. 본 \S 3
      에서는 형태에 대한 가정만 두며, 정량은 \S\ref{sec:rho_eq} 에서.
\end{enumerate}

\subsection{연속 모델의 물리적 근거}
\label{sec:physical_grounding}

흑연 음극의 Li 인터칼레이션을 연속 화학 퍼텐셜 분포로 모델링하는 접근은 다음 두
선행 연구에서 정당화된다.

\begin{itemize}
\item \textbf{Cahn-Hilliard reaction dynamics for graphite} (Guo, Bazant, et al.,
      \emph{J. Phys. Chem. Lett.} \textbf{7}, 2151 (2016)) — 흑연 인터칼레이션 시
      Li 화학 퍼텐셜이 공간적으로 연속이며, 상변이가 phase coexistence + chemical
      potential gradient 의 형태로 진행됨을 직접 광학 이미징으로 보임.
\item \textbf{Free energy model coupled to graphene stacking}
      (Rykner, Chandesris, \emph{J. Phys. Chem. C} \textbf{126}, 11425 (2022)) —
      자유 에너지가 stacking sequence 와 결합되어 chemical potential 의 연속
      변화를 만들며, 단순 $N_p$ 이산 staging 으로 분해되지 않음을 보임.
\end{itemize}

두 연구는 본 Chapter 1 의 연속 spine 의 \textbf{물리적 근거} 이며, ver5 의 이산
$N_p$ 분해가 부분적 근사일 수 있음을 시사한다.

\subsection{연속 → 이산 일방 유도}
\label{sec:continuous_to_discrete}

본 spine 은 \textbf{연속 분포 $\rhomu$ 를 primary state 로 두지만}, 기존
\texttt{ver5} 와 \texttt{ver1\_rechecked2} 에서 사용되어 온 이산 $\xi_j$ 양과의
호환을 위해 다음 \emph{일방 유도} 를 명시한다.

각 peak window $j$ 를 화학 퍼텐셜 구간 $[\mucont_{j,\mathrm{lo}}(T),\,
\mucont_{j,\mathrm{hi}}(T)]$ 로 정의 하면 ($j = 1, 2, \ldots, N_p$, $N_p$ 는
관측에 따라 결정되는 정성적 분류 수), 해당 peak 의 누적 Li 전하량은
\begin{equation}
Q_j(q, T)
\;=\;
\int_{\mucont_{j,\mathrm{lo}}(T)}^{\mucont_{j,\mathrm{hi}}(T)} \dd\mucont \;
\rhomu \cdot \nmu
\label{eq:peak_charge}
\end{equation}
이고, 진행률은
\begin{equation}
\xi_j(q, T)
\;=\;
\frac{Q_j(q, T) - Q_j(q_0, T)}{Q_{j,\tot}(T)},
\qquad
Q_{j,\tot}(T) = Q_j(q = 1, T) - Q_j(q = 0, T)
\label{eq:xi_from_rho}
\end{equation}
로 정의된다 ($q_0$ 은 진행률의 기준 시점, 보통 $q_0 = 0$).

\begin{note}
식 \eqref{eq:xi_from_rho} 는 \emph{일방 유도} 이다. 즉 $\rhomu$ 가 주어지면 $\xi_j$
가 결정되지만, 역방향 (이산 $\xi_j$ 만 주어졌을 때 연속 $\rhomu$ 를 복원) 은 본
spine 에서 사용되지 않는다. 역방향 복원은 정보 손실을 동반하며, Charter \S 4 FB8
(silent $\mucont$ 이산화) 에 해당하므로 정의식에서 배제된다.

따라서 본 Chapter 1 에서 $\xi_j$ 가 등장하는 모든 위치 (\S 13 피팅 절차 등) 는
연속 $\rhomu$ 모델의 결과를 이산 형태로 \emph{보고} 하는 단계일 뿐, 모델 자체의
state 가 아니다.
\end{note}

\subsection{peak window 의 운용적 정의}
\label{sec:peak_window_operational}

식 \eqref{eq:peak_charge}--\eqref{eq:xi_from_rho} 에서 peak window
$[\mucont_{j,\mathrm{lo}}, \mucont_{j,\mathrm{hi}}]$ 의 정의는 \emph{운용적
(operational)} 이다. 즉, $\rhomu$ 의 형태가 결정된 뒤에 peak 의 valley 사이
중간점, FWHM-기반 영역, 또는 모멘트-기반 영역 등으로 정의할 수 있다. 구체적
선택은 \S\ref{sec:fitting_procedure} (Phase E10 피팅 절차) 에서 다룬다.

\textbf{본 \S 3 의 핵심 결론.} 본 Chapter 1 의 primary state 는 연속 분포 $\rhomu$
이며, 이산 $\xi_j$ 와 그에 부속된 $U_j, w_j, k_j, Q_{j,\tot}$ 등은 모두 식
\eqref{eq:xi_from_rho} 와 그 변형으로 \emph{유도되는} 양이다. 이 일방향 정의가
사용자가 지적한 ``적분을 모 아니면 도 스텝 펑션 형태로 가정'' 한 결함의 근본 회피
점이다.

% ====================================================================
\section{음극 전위의 세 가지 구분: $V_n$, $V_{n,\app}$, $V_{n,\drive}$}
\label{sec:three_potentials}
% ====================================================================

본 Chapter 1 은 \texttt{ver1\_rechecked2} \S 6.1 (line 204--238) 의 세 가지 전위
구분을 정신적으로 보존하면서, 새 연속 $\rhomu$ spine 위에 적합하게 재정의한다.

\subsection{$V_n$: 음극 내부 전위 (Eq.\,\ref{eq:spine_summary} (2) 의 implicit 해)}
\label{sec:V_n_internal}

음극 내부 전위 $V_n$ 은 주어진 $(q, T, \rhomu)$ 에서 식 \eqref{eq:spine_summary} (2)
의 \emph{implicit 해} 로 정의된다. 즉,

\begin{equation}
V_n
\;=\;
\bigl\{
V_n \in \mathbb{R} \;\big|\;
Q_{\cell}\, q
\;=\;
Q_{\bg}(V_n, T)
\;+\;
\int_{\mucont_{\min}}^{\mucont_{\max}} \dd\mucont \;
\rhomu \cdot \nmu
\bigr\}.
\label{eq:V_n_implicit}
\end{equation}

$V_n$ 은 외부에서 lookup 으로 주어지는 함수가 아니다 (ver5 의 $V_{n,\OCV}(q, T)$
방식 폐기). 또한 $V_n$ 은 단일 boundary 값이고, $\rhomu$ 의 분포 내부 좌표
$\mucont$ 와 \emph{개념적으로 구분} 된다 (자세한 대응 관계는
\S\ref{sec:mu_V_n_correspondence}).

식 \eqref{eq:V_n_implicit} 의 root-find 는 매 timestep 에서 수치적으로 수행되며,
이를 \textbf{Loop A} (Phase E1 \S D 분류) 라 부른다. Loop A 의 해 존재 조건과
수치 안정성 조건은 \S\ref{sec:charge_balance} (Phase E4) 에서 다룬다.

\subsection{$V_{n,\app}$: 관측 전위}
\label{sec:V_n_app}

실험에서 측정되는 apparent 음극 전위는 식 \eqref{eq:Vapp_def} 으로 이미 정의되어
있다. 본 \S 4.2 는 그 식의 역할을 재확인한다.

\begin{equation}
V_{n,\app}(q, T, |I|)
\;=\;
V_n(q, T, \rhomu) \;+\; s_I |I|\, R_n(q, T, |I|).
\label{eq:V_n_app_role}
\end{equation}

여기서 두 가지가 새롭다.

\begin{itemize}
\item $V_n$ 의 의존성에 $\rhomu$ 가 명시적으로 등장한다 (Eq.\,\ref{eq:V_n_implicit}
      의 해 의존성). ver5 와 달리 $V_n$ 자체가 분포 $\rhomu$ 의 함수이다.
\item $R_n$ 의 $|I|$ 의존성은 \texttt{ver1\_rechecked2} 에서 이미 도입된 형태로
      유지된다. ver5 의 $R_n(q, T)$ 보다 일반적이며, 분극 효과의 비선형성을
      포함한다.
\end{itemize}

\subsection{$V_{n,\drive}$: 반응 커널의 구동 전위}
\label{sec:V_n_drive}

연속 반응성 커널 $\Sreact$ 의 구동력 계산에 사용되는 전위는 $V_{n,\drive}$ 로
표기한다. Reduced model 에서는

\begin{equation}
V_{n,\drive} \;\approx\; V_{n,\app}
\label{eq:V_n_drive_reduced}
\end{equation}

근사를 사용한다. 이 근사는 다음 조건이 동시에 충족될 때 유효하다.

\begin{enumerate}
\item $R_n$ 에 전하 이동 분극 성분이 \emph{이미 포함} 되어 있으며, $\Sreact$ 가
      그 성분을 \emph{다시 포함하지 않는다} (중복 반영 금지).
\item 외부 회로의 전류 분배가 단일 음극 반응 채널로 통합되어 다룰 수 있다
      (다중 채널 분리 표현은 Phase E6 reaction-resolved 확장에서 다룸).
\end{enumerate}

Reaction-resolved 확장형 ($V_{n,\drive}$ 가 Butler-Volmer 과전압 또는 다른
구체 형태로 분리 표현되는 경우) 은 Phase E6 (\S\ref{sec:continuous_reactivity}
신단원) 에서 정의되며, 본 \S 4.3 에서는 reduced approximation 만 명시한다.

\subsection{세 전위의 역할 차이}
\label{sec:three_potentials_roles}

\begin{mdframed}
\textbf{역할 요약.}

\begin{tabular}{p{0.18\textwidth} p{0.74\textwidth}}
\toprule
전위 & 역할 \\
\midrule
$V_n$ & 모델의 내부 상태. 식 \eqref{eq:V_n_implicit} 의 implicit 해. $\rhomu$ 와
        $Q_{\bg}$ 와 결합되어 정의됨. \\
$V_{n,\app}$ & 외부 관측량. $V_n$ 과 옴성 분극 항의 합. ICA/DVA 관측식
        (\S\ref{sec:ICA_DVA}) 의 미분 좌표. \\
$V_{n,\drive}$ & 반응성 커널 $\Sreact$ 의 구동 좌표. Reduced 에서 $V_{n,\app}$
        근사. \\
\bottomrule
\end{tabular}
\end{mdframed}

세 전위를 혼동하지 않는 것은 본 Chapter 1 의 자기 일관성 (Charter \S 6 audit
Dim \#11 의 system-fidelity 검증) 의 핵심이다. 예를 들어 평형 분포 $\rhoeq$ 는
$V_n$ 의 함수로 정의되어야 하며 (반응의 구동력이 아닌 평형 위치), 반응성 커널
$\Sreact$ 는 $V_{n,\drive}$ 의 함수로 정의되어야 한다 (Phase E6).

\subsection{$\mucont$ 와 $V_n$ 의 boundary 대응}
\label{sec:mu_V_n_correspondence}

연속 화학 퍼텐셜 좌표 $\mucont$ 는 \emph{분포 $\rhomu$ 의 내부 좌표} 이고,
$V_n$ 은 분포 전체에 대해 정의되는 \emph{단일 boundary 값} 이다. 두 양은 다음
의미에서 연결된다.

평형 ($\rhomu = \rhoeq$) 조건에서 본 spine 의 charge balance
(Eq.\,\ref{eq:V_n_implicit}) 를 만족하는 $V_n$ 의 값은 $\rhoeq$ 의 분포가 최대로
점유된 $\mucont$ 영역의 \emph{대표값} (boundary, 또는 chemical potential 의 평형
조건의 결과) 으로 해석된다. 비평형 ($\rhomu \neq \rhoeq$) 조건에서도 동일한 식
이 $V_n$ 을 정의한다.

\begin{mdframed}
\textbf{새 boundary 대응 = 기존 $V_{n,\OCV}$ 외부 함수의 대체.}

\texttt{ver5} 의 $V_{n,\OCV}(q, T)$ 는 외부 lookup 함수로 가정되었고, 평형 진행률
$\xi_{j,\eq}$ 의 정의에 직접 인입되었다. 이는 charge balance 가 모델 외부에서
이미 풀린 것으로 가정한 셈이다.

본 Chapter 1 v0.1 에서는 $V_n$ 이 식 \eqref{eq:V_n_implicit} 의 implicit 해로
\emph{내부적으로 결정} 되며, 평형 시 그 값이 자연스럽게 $V_{n,\OCV}(q, T)$ 의 역할
을 한다 (derived quantity, Phase E1 \S C 참조). 따라서 $V_{n,\OCV}$ 는 더 이상
모델의 입력이 아니라, 모델이 평형 조건 하에서 \emph{만들어내는} 양이다.
\end{mdframed}

\subsection{본 \S 의 \S 1 spine 과의 연결 및 \S 5 진입}
\label{sec:section_4_to_5_bridge}

본 \S 4 는 \S 1 식 \eqref{eq:spine_summary} 의 성분 (3) ($V_n$ 의 implicit 해 정의)
와 성분 (4) ($V_{n,\app}$ 정의) 를 본문화하였다. 다음 \S 5 (charge balance 중심식,
Phase E4 작성) 에서는 성분 (2) 의 정식 유도와 해 존재 조건, 수치 안정성, 총용량
정합, 단조성 제약을 다룬다.

\end{document}
